\scanpage{237}\mspage{92}
$W$-coasphérique (donc $\in W$) pour tt \deleted{\unsure{fonct}} $W$
satisfaisant l'axiome $\mathrm{Loc}\ 8$ des limites
inductives. En particulier, elle est $W_0$-\unsure{coasph},
i.e. un foncteur cofiltrant \deleted{\unsure{ou}} (et $=$
seulmt fbt cofiltrant) est \deleted{\illegible{1}} \uline{cofinal}.

La raison principale pour m'intéresser
aux flèches cofiltrantes est contenue
dans la

\uline{Proposition} 7.4. Soit $f:X\longrightarrow Y$ \inserted{dans $\operatorname{Cat}$,} loc. cofiltrant
(resp. loc. \uline{fbt.} cofiltrant)
Alors : \deleted{pour toute}

a) Pour tout $X'$ étale sur $X$, le composé $f':$
$X'\unsure{-|\to}X-Y$ est loc. cofiltrant\inserted{ (resp. \ldots)}, et pour
toute factorisation de $f'$ en $X'\xrightarrow{g'}Y'\longrightarrow Y$,
avec $Y'-Y$ étale, $g'$ est loc cofiltrant (resp. \ldots).

b) \deleted{\unsure{Les parties libres}}, \inserted{$\forall x$ dans $X$,} le \unsure{morphisme} sur les localisés
\deleted{\unsure{sont}} $f_{/x}:X_{/x}\unsure{-|\to}Y_{/y}\qquad(y=f(x))$
est \uline{cofiltrant} (resp. fbt cofiltrant),
\deleted{donc $\unsure{W_0}$-coasphérique, i.e. \uline{cofinal}.}

\deleted{Pour} c) Si $x$, \deleted{\unsure{est un}} comme dans b), et
$U$ est un ouvert de $X_{/x}$, désignant
par $V$ l'ouvert de $Y_{/y}$ engendré par
l'image de $f_{/x}$ (i.e. l'ouvert formé des
$z$ dans $Y_{/y}$ tels que $\,{}_{z\backslash}(X_{/x})$ soit non vide),
la flèche induite $g:U\longrightarrow V$ \unsure{est|\deleted{est}} \inserted{cofiltrante (resp. fbt \unsure{cofiltrte}), \unsure{a fortiori}} \uline{cofinale},
(et $\unsure{\overline{m}|=}$ $W$-coasphérique, pour toute $W$
satisfaisant l'axiome des limites $\mathrm{Loc}\ 8$).
% slanted
\margin{Il suffit ici qu'on soit dans l'hyp. \unsure{resp}\textsuperscript{\unsure{ve}} \deleted{plus} ($f$ \uline{fbt}. loc. cofilt\unsure{(ant)}) \unsure{Sauf} pour le \unsure{NB ici}.}
% archive sheet number: 237
